\documentclass[9pt,xcolor=svgnames]{beamer} %Beamer

\usepackage{palatino} %font type
\usepackage{svg}
\usepackage{stackengine}
\usefonttheme{metropolis} %Type of slides
\usefonttheme[onlymath]{serif} %font type Mathematical expressions
\usetheme[progressbar=frametitle,titleformat frame=smallcaps,numbering=counter,sectionpage=progressbar]{metropolis} %This adds a bar at the beginning of each section.
\useoutertheme[subsection=false]{miniframes} %Circles in the top of each frame, showing the slide of each section you are at



\usepackage{appendixnumberbeamer} %enumerate each slide without counting the appendix
\setbeamercolor{progress bar}{fg=Maroon!70!Coral} %These are the colours of the progress bar. Notice that the names used are the svgnames
\setbeamercolor{title separator}{fg=DarkSalmon} %This is the line colour in the title slide
\setbeamercolor{structure}{fg=black} %Colour of the text of structure, numbers, items, blah. Not the big text.
\setbeamercolor{normal text}{fg=black!87} %Colour of normal text
\setbeamercolor{alerted text}{fg=DarkRed!60!Gainsboro} %Color of the alert box
\setbeamercolor{example text}{fg=Maroon!70!Coral} %Colour of the Example block text

\definecolor{mydarkred}{HTML}{A62618}

\definecolor{highlight1}{HTML}{FFE070}
\definecolor{highlight2}{HTML}{D0D0FF}

\setbeamercolor{palette primary}{bg=NavyBlue!50!DarkOliveGreen, fg=white} %These are the colours of the background. Being this the main combination and so one. 
\setbeamercolor{palette secondary}{bg=NavyBlue!50!DarkOliveGreen, fg=white}
\setbeamercolor{palette tertiary}{bg=NavyBlue!40!Black, fg= white}
\setbeamercolor{section in toc}{fg=NavyBlue!40!Black} %Color of the text in the table of contents (toc)

\setbeamercovered{transparent}

\usepackage{ulem}
\usepackage{arevmath}
\usepackage{fontenc}

\usepackage{pgfpages}
% \setbeameroption{show notes}
% \setbeameroption{show notes on second screen=right}

\newcommand{\xmark}{\scalebox{0.75}{$\ballotx$}}
%These next packages are the useful for Physics in general, you can add the extras here. 
\usepackage{amsmath,amssymb}
\usepackage{slashed}
\usepackage{cite}
\usepackage{relsize}
\usepackage{caption}
\usepackage{subcaption}
\usepackage{multicol}
\usepackage{booktabs}
\usepackage[scale=2]{ccicons}
\usepackage{pgfplots}
\usepgfplotslibrary{dateplot}
\usepackage{geometry}
\usepackage{xspace}
\newcommand{\themename}{\textbf{\textsc{bluetemp}\xspace}}%metropolis}}\xspace}
\usepackage{nameref}
\makeatletter
\newcommand*{\currentname}{\@currentlabelname}
\makeatother

% \usepackage{soul}
% \input{highlight.sty}

% \usepackage{tcolorbox}
% \tcbuselibrary{skins, theorems}
% \tcbset{fonttitle=\footnotesize}
% \newcommand{\highlight}[2][yellow]{
%   \tcboxmath[colback=#1]{#2}
% }

\usepackage{graphicx}
\usepackage{pgf}
\usepackage{lipsum}
\usepackage{hyperref}
\usepackage{svg}
\usepackage{float}
\usepackage{aliascnt}

\newfloat{eqfloat}{h}{eqflts}
\floatname{eqfloat}{Equation}

\newcommand{\vspacer}{\vphantom{\left(\dfrac{\sum_1^1 f \left(\sum_1^1\right)}{\sum_1^1 f \left(\sum_1^1\right)}\right)}}

\newcommand{\highlight}[2][yellow]{\mathchoice%
  {\colorbox{#1}{$ \vspacer \displaystyle#2$}}%
  {\colorbox{#1}{$ \vspacer \textstyle#2$}}%
  {\colorbox{#1}{$ \vspacer \scriptstyle#2$}}%
  {\colorbox{#1}{$ \vspacer \scriptscriptstyle#2$}}}%


\newcommand{\mdr}[1]{\textbf{\color{mydarkred} #1} }
\newcommand{\RomanNumeralCaps}[1]
    {\MakeUppercase{\romannumeral #1}}

\graphicspath{ {./img/} }

\title{{\huge Project Presentation}}
\subtitle{Group 7}
\author[Name]{
Stefano Pea \#561020 \newline Simone Ianniciello \#581201 \newline Claudia Gentili \#581522
} %With inst, you can change the institution they belong
\institute[uni]{Department of Computer Science \\ University of Pisa}
\date{\textbf {January 25 2024}} %Here you can change the date
\usepackage{tikz}
\usetikzlibrary{calc}
\titlegraphic { 
\begin{tikzpicture}[overlay,remember picture]
    \node[left=-0.2cm] at (current page.23){
        \includegraphics[width=4.5cm]{img/brushed_unipi.png}
    };
    \node[left=0.8cm] at (current page.29){
        %\includegraphics[width=10.6cm]{img/ferragina.png}
    };
    % \node[text width=5cm, right=-5cm, ] at (current page.-30){
    %     \color{White}
    %     {DataMining Paper Presentation}
    % };
\end{tikzpicture}
}
\pgfmathsetseed{42}
\newcommand{\scribblecancel}[1]{%
  \tikz[baseline=(X.base)]{
    \node[inner sep=0pt, outer sep=0pt] (X) {#1};
    \coordinate (start) at (X.south west);
    \foreach \x in {1,...,40} {
      \pgfmathsetmacro{\randx}{random(0,100)/100}
      \pgfmathsetmacro{\randy}{random(0,100)/100}
      \coordinate (current) at ($(X.south west)!\randx!(X.south east) + \randy*(0,0.4)$);
      \draw[line cap=round, line join=round, line width=.5pt] (start) -- (current);
      \coordinate (start) at (current);
    }
  }%
}

% \titlegraphic{\vspace{-0.9cm}\hfill\includegraphics[scale=0.23]{brushed_unipi.png}} %You can modify the location of the logo by changing the command \vspace{}. 

\usepackage{wrapfig}

\begin{document}
{
% #003c71
\setbeamercolor{background canvas}{bg=NavyBlue!50!DarkOliveGreen, fg=White}
\setbeamercolor{normal text}{fg=White}
\maketitle
}%This is the colour of the first slide. bg= background and fg=foreground

\metroset{titleformat frame=smallcaps} %This changes the titles for small caps

\section{Data Preparation}




\begin{frame}{Data Understanding\/Preparation}
    

    \begin{exampleblock}{Initial Preparation}
       
        \begin{itemize}
            \item Merged incidents, poverty and state district house datasets with population dataset
    
        \end{itemize}
        
    \end{exampleblock}


    \begin{exampleblock}{Problems}
        \begin{itemize}
            \item Out of range values
            \item Missing congressional districts
            \item Inconsistencies in the participants groups
            \item Missing values
        \end{itemize}
    
    \end{exampleblock}


\end{frame}



%%%%%%%%%%%%%%%%%%%%%%%%%%%%%%%%%%%%%%%%%%%%%%%%%%%%%%%%%%%%%%%%%%%%%%%%%%


\begin{frame}
    \frametitle{Cleaned Data Analysis I}
    \begin{figure}
        \centering
        \includegraphics[width=.99\textwidth]{img/DU/dist.png}
        \label{corr}
    \end{figure}
\end{frame}


%%%%%%%%%%%%%%%%%%%%%%%%%%%%%%%%%%%%%%%%%%%%%%%%%%%%%%%%%%%%%%%%%%%%%%%%%%


\begin{frame}
    \frametitle{Cleaned Data Analysis II}
    \begin{figure}
        \centering
        \includegraphics[width=.75\textwidth]{img/DU/corrScaled.png}
        \label{corr}
    \end{figure}
\end{frame}

%%%%%%%%%%%%%%%%%%%%%%%%%%%%%%%%%%%%%%%%%%%%%%%%%%%%%%%%%%%%%%%%%%%%%%%%%%

\begin{frame}
    \frametitle{Indicators}
    \begin{figure}
        \centering
        \includegraphics[width=.8\textwidth]{img/DU/distInd.png}
        \label{corr}
    \end{figure}
\end{frame}
\section{Clustering}

\begin{frame}
    \frametitle{Clustering}
    \begin{exampleblock}{Preparation }
        
    
        \begin{itemize}
            \item IQR method for outliers detection
            \item Min-Max scaling
            \item Silhouette score for evaluation
        \end{itemize}
    \end{exampleblock}    
\end{frame}

%%%%%%%%%%%%%%%%%%%%%%%%%%%%%%%%%%%%%%%%%%%%%%%%%%%%%%%%%%%%%%%%%%%%%%%%%%

\begin{frame}
    \frametitle{Clustering}

\begin{exampleblock}{K-Means}
    \begin{itemize}
        \item Entire dataset
        \item Parameter: $k$ (number of clusters)
    \end{itemize}
\end{exampleblock}
\begin{exampleblock}{DBSCAN}
    \begin{itemize}
        \item State: California
        \item Parameters: {
            \begin{itemize}
                \item $\epsilon$ (neighborhood radius)
                \item $n$ (core points neighborhood minimum size)
            \end{itemize}
        }
        
    \end{itemize}
\end{exampleblock}
\begin{exampleblock}{Hierarchical}
    \begin{itemize}
        \item State: California
        \item Parameter: inter-cluster distance
    \end{itemize}
\end{exampleblock}

\end{frame}

%%%%%%%%%%%%%%%%%%%%%%%%%%%%%%%%%%%%%%%%%%%%%%%%%%%%%%%%%%%%%%%%%%%%%%%%%%

\begin{frame}{K-means}
    \begin{exampleblock}{Selected features:}
        \begin{itemize}
            \item state\_population
            \item timestamp
            \item p\_injured
            \item month\_cd\_change\_min\_age\_participants
            \item month\_cd\_ratio\_males
        \end{itemize}        
    \end{exampleblock}


    
    \begin{table}[]
        \centering
        \begin{tabular}{|lc|}
    
            \hline
            \textbf{Clustering}  & C6 \\
            \textbf{K} & 2\\
            \textbf{Clusters sizes} & (59466, 35725) \\
            \textbf{Silhouette mean} & 0.453778 \\
            \textbf{Silhouette variance} & 3.048582e-08 \\
            \hline

        \end{tabular}
        \label{tab:my_label}
    \end{table}
\end{frame}

%%%%%%%%%%%%%%%%%%%%%%%%%%%%%%%%%%%%%%%%%%%%%%%%%%%%%%%%%%%%%%%%%%%%%%%%%%

\begin{frame}{K-means}
    \begin{figure}
        \centering
        \includegraphics[width=.99\textwidth]{img/clustering/kmeans.png}
        \label{kmeans}
    \end{figure}
\end{frame}


%%%%%%%%%%%%%%%%%%%%%%%%%%%%%%%%%%%%%%%%%%%%%%%%%%%%%%%%%%%%%%%%%%%%%%%%%%

\begin{frame}{DBSCAN}

    \begin{exampleblock}{Selected features:}
        \begin{itemize}
            \item latitude
            \item min\_age\_participants
            \item n\_killed
        \end{itemize}
           
    \end{exampleblock}
    
    \begin{table}[]
        \centering
        \begin{tabular}{|lc|}
            \hline
              \textbf{Clustering}  & C4 \\
              $\epsilon$  & 0.6 \\
              $n$ & 21 \\
              \textbf{K} & 2\\
              \textbf{Clusters sizes} & (2088, 3938) \\
              \textbf{Silhouette} & 0.670266 \\
              \hline
        \end{tabular}
        \label{tab:my_label}
    \end{table}
\end{frame}

%%%%%%%%%%%%%%%%%%%%%%%%%%%%%%%%%%%%%%%%%%%%%%%%%%%%%%%%%%%%%%%%%%%%%%%%%%

\begin{frame}{DBSCAN}
    \begin{figure}
    \centering
    \includegraphics[width=.99\textwidth]{img/clustering/dbscan.png}
    \label{dbscan}
\end{figure}
\end{frame}

%%%%%%%%%%%%%%%%%%%%%%%%%%%%%%%%%%%%%%%%%%%%%%%%%%%%%%%%%%%%%%%%%%%%%%%%%%

\begin{frame}{Hierarchical Clustering}

    \begin{exampleblock}{Selected features:}
        \begin{itemize}
            \item latitude
            \item min\_age\_participants
            \item n\_killed
        \end{itemize}
    \end{exampleblock}

    
    
    \begin{table}[]
        \centering
        \begin{tabular}{|lc|}
            \hline  
            \textbf{Clustering}  & C4 \\
              \textbf{Best Method} & all\\
              \textbf{K} & 2\\
              \textbf{Clusters sizes} & (2088, 3938) \\
              \textbf{Silhouette} & 0.670266 \\
              \hline
        \end{tabular}
        \label{tab:my_label}
    \end{table}
\end{frame}

%%%%%%%%%%%%%%%%%%%%%%%%%%%%%%%%%%%%%%%%%%%%%%%%%%%%%%%%%%%%%%%%%%%%%%%%%%

\begin{frame}{Hierarchical Clustering}
    \begin{figure}
    \centering
    \includegraphics[width=.99\textwidth]{img/clustering/hier.png}
    \label{hier}
\end{figure}
\end{frame}

%%%%%%%%%%%%%%%%%%%%%%%%%%%%%%%%%%%%%%%%%%%%%%%%%%%%%%%%%%%%%%%%%%%%%%%%%%

\begin{frame}{DBSCAN and Hierarchical Clusters}
    \begin{figure}
    \centering
    \includegraphics[scale=0.29]{img/clustering/dbscan1.png}
    \includegraphics[scale=0.31]{img/clustering/hier2.png}
    \label{hier}
\end{figure}
\end{frame}
\section{Classification}

\begin{frame}[standout]{\currentname: Problem Description}
    Goal: predict the presence of a fatality in a gun incident.
\end{frame}

\begin{frame}<1-5>{\currentname: Feature Selection}
    \begin{exampleblock}{\textsc{Obvious Features to Remove}}
        All features directly related to the target variable.
    \end{exampleblock}
    \begin{exampleblock}<2->{\textsc{Less Obvious}}
        \begin{itemize}
            \item n\_killed = {\only<3->{\color{red}}n\_participants} - n\_injured - n\_unharmed - n\_arrested
            \item n\_participants = {\only<4->{\color{red}}n\_males} + n\_females
            \item n\_participants = {\only<4->{\color{red}}n\_adult} + n\_teen + n\_child
        \end{itemize}
    \end{exampleblock}
    \begin{exampleblock}<5->{\textsc{Lesser Obvious Problem}}
        $\sim50\%$ of datapoints have n\_participants = 1.
    \end{exampleblock}
\end{frame}

\begin{frame}{\currentname: Data Preparation}
    \begin{exampleblock}{\textsc{Scaling}}
        \vspace*{0.5em}
        StandardScaler from sklearn.
    \end{exampleblock}
    \begin{exampleblock}{\textsc{Train-Valid-Test Splits}}
        \vspace*{-1em}
        \begin{flalign*}
            \text{Train} &= 70\%& \\
            \text{Valid} &= 15\%& \\
            \text{Test}  &= 15\%&
        \end{flalign*}
    \end{exampleblock}
    \vspace*{-1em}
    \begin{exampleblock}{\textsc{Balancing}}
        \vspace*{0.5em}
        $75\%$ of 0-labelled datapoints. \\
        \vspace*{0.5em}
        \begin{minipage}{0.47\textwidth}
            Training balanced with SMOTE.
        \end{minipage}
        \begin{minipage}{0.47\textwidth}
            Validation and test balanced with random sampling with replacement.
        \end{minipage}
    \end{exampleblock}
\end{frame}

\begin{frame}{\currentname: Models}
    \begin{exampleblock}{Decision Tree}
        \begin{itemize}
            \item[] 16 levels
            \item[] Gini criterion
            \item[] 2 minimum samples per leaf
            \item[] 1 minimum sample per split
        \end{itemize}
    \end{exampleblock}
    \begin{exampleblock}{Neural Network}
        \begin{itemize}
            \item[] 4 hidden layers
            \item[] [32, 16, 8, 4] neurons per layer
            \item[] Sigmoid activation
            \item[] 0.1 dropout
        \end{itemize}
    \end{exampleblock}
\end{frame}

\begin{frame}{\currentname: Decision Tree Results I}
    \begin{figure}
        \centering
        \includegraphics[width=.99\textwidth]{CS/decision_tree.png}
        \caption*{First 4 levels of the decision tree.}
    \end{figure}
\end{frame}

\begin{frame}{\currentname: Decision Tree Results II}
    \begin{figure}
        \centering
        \includegraphics[width=.99\textwidth]{CS/decision_tree_cm.png}
        \caption*{Confusion matrix of the decision tree in TR/VL/TE splits.}
    \end{figure}
    \begin{table}[h]
        \centering
        \begin{tabular}{l|c|c|c|}
            Split & Train & Valid & Test \\
            \hline
            Accuracy  & 0.93 & 0.77 & 0.76 \\
            Precision & 0.96 & 0.94 & 0.94 \\
            Recall    & 0.90 & 0.70 & 0.70 \\
            F1        & 0.93 & 0.80 & 0.77 \\        
        \end{tabular}
    \end{table}
\end{frame}

\begin{frame}{\currentname: Neural Network Results}
    \begin{figure}
        \centering
        \includegraphics[width=.99\textwidth]{CS/nn_cm.png}
        \caption*{Confusion matrix of the Neural Network in TR/VL/TE splits.}
    \end{figure}    
    \begin{table}[h]
        \centering
        \begin{tabular}{l|c|c|c|}
            Split & Train & Valid & Test \\
            \hline
            Accuracy  & 0.83 & 0.80 & 0.79 \\
            Precision & 0.87 & 0.87 & 0.87 \\
            Recall    & 0.80 & 0.76 & 0.75 \\
            F1        & 0.83 & 0.81 & 0.81 \\        
        \end{tabular}
    \end{table}
\end{frame}

\begin{frame}{\currentname: Models Comparison}
    \begin{figure}
        \centering
        \includegraphics[width=.99\textwidth]{CS/roc-prc.png}
        \caption*{ROC/PRC curves of the two models in TR/VL/TE splits.}
    \end{figure}
\end{frame}
\section{Time Series}





\begin{frame}<1-2>{\currentname: City Score Definition}

    

    \begin{exampleblock}{\textsc{Score = }}
        \begin{flalign*}
            & {\only<2->{\color{red}} n\_killed\_week }*  {\only<2->{\color{green}} 1} + {\only<2->{\color{red}} n\_injured\_week} *  {\only<2->{\color{green}} 0.8} + {\only<2->{\color{red}} n\_arrested\_week} * {\only<2->{\color{green}} 0.5} \\
            &+ {\only<2->{\color{red}} n\_child\_week}  * {\only<2->{\color{green}} 1} + {\only<2->{\color{red}} n\_teen\_week }* {\only<2->{\color{green}} 0.8} + {\only<2->{\color{red}} n\_adult\_week} * {\only<2->{\color{green}} 0.5} \\
        \end{flalign*}
        \begin{flalign*}
            &n\_killed\_week = \text{number of killed people in a specific week} \\
            &n\_injured\_week = \text{number of injured people in a specific week} \\
            &n\_arrested\_week = \text{number of arrested people in a specific week} \\
            &n\_child\_week = \text{number of child participants in a specific week} \\
            &n\_teen\_week = \text{number of teen participants in a specific week} \\
            &n\_adult\_week = \text{number of adult participants in a specific week} \\
        \end{flalign*}

    \end{exampleblock}
\end{frame}

%%%%%%%%%%%%%%%%%%%%%%%%%%%%%%%%%%%%%%%%%%%%%%%%%%%%%%%%%%%%%%%%%%%%%%%%%%

\begin{frame}{\currentname: Time Series Overview}

    \begin{figure}
        \centering
        \includegraphics[width=.99\textwidth]{img/ts/a.png}
        \caption{Generated Time Series}
        \label{GTS}
    \end{figure}
    

\end{frame}


%%%%%%%%%%%%%%%%%%%%%%%%%%%%%%%%%%%%%%%%%%%%%%%%%%%%%%%%%%%%%%%%%%%%%%%%%


%TODO Make the various parts appear in order
\begin{frame}<1-3>{Clustering}
    \begin{minipage}{0.43\textwidth}
    
        
        \begin{exampleblock}{\textsc{Types of Clustering:}}
            \vspace{.8em}

            \begin{itemize}
                \item Standard Clustering
                \item Feature Clustering
            \end{itemize}
            \end{exampleblock}        
        
        
            \begin{exampleblock}<3->{\textsc{Algorithms utilized:}}        
                    
                \begin{itemize}
                
                    \item K-Means on Time series
                \end{itemize}
            \end{exampleblock}
        
    \end{minipage}
    \begin{minipage}{0.43\textwidth}
        

        \visible<2->{
        \begin{tabular}{|l|l|}
            \hline
            \textbf{Feature} & \textbf{Description} \\
            \hline
            'avg' & Mean (average) of the values. \\
            'std' & Standard deviation of the values. \\
            'var' & Variance of the values. \\
            'med' & Median of the values. \\
            '10p' & 10th percentile of the values. \\
            '25p' & 25th percentile of the values. \\
            '50p' & 50th percentile of the values. \\
            '75p' & 75th percentile of the values. \\
            '90p' & 90th percentile of the values. \\
            'iqr' & Interquartile range. \\
            'cov' & Coefficient of variation. \\
            'skw' & Skewness of the values. \\
            'kur' & Kurtosis of the values. \\
            \hline
        \end{tabular}
        }
        
    \end{minipage}
    
\end{frame}

%%%%%%%%%%%%%%%%%%%%%%%%%%%%%%%%%%%%%%%%%%%%%%%%%%%%%%%%%%%%%%%%%%%%%%%%%%

\begin{frame}{Silhouette Score}

    \begin{table}[htbp]
        \centering
        \begin{tabular}{|c|c|c|}
          \hline
          \textbf{n\_clusters} & \textbf{silhouette score original} & \textbf{silhouette score features}\\
          \hline
          2 & 0.812281 & 0.984422 \\
          3 & 0.781509 & 0.936071 \\
          4 & 0.625345 & 0.872255 \\
          5 & 0.458336 & 0.825605 \\
          6 & 0.458559 & 0.765932 \\
          7 & 0.461002 & 0.765437 \\
          8 & 0.424799 & 0.689629 \\
          9 & 0.413574 & 0.690026 \\
          \hline
        \end{tabular}
        \label{Silhouette Score for the different Scaling}
        \caption{Silhouette Score the Original Time Series and their Features}
      \end{table}      

\end{frame}
        


%%%%%%%%%%%%%%%%%%%%%%%%%%%%%%%%%%%%%%%%%%%%%%%%%%%%%%%%%%%%%%%%%%%%%%%%%%

\begin{frame}{Motifs and Anomaly Detection}

    \begin{figure}
        \centering
        \includegraphics[width=.99\textwidth]{img/ts/Motifs.png}
      
        \label{MOT}
    \end{figure}
    
    
        
\end{frame}
        

%%%%%%%%%%%%%%%%%%%%%%%%%%%%%%%%%%%%%%%%%%%%%%%%%%%%%%%%%%%%%%%%%%%%%%%%%%

\begin{frame}{Motifs and Anomaly Detection 2}

    \begin{figure}
        \centering
        \includegraphics[width=.8\textwidth]{img/ts/ChicagoMotifs.png}
      
        \label{CHGMOT}
    \end{figure}
    
    
        
\end{frame}


%%%%%%%%%%%%%%%%%%%%%%%%%%%%%%%%%%%%%%%%%%%%%%%%%%%%%%%%%%%%%%%%%%%%%%%%%%


\begin{frame}{Motifs and Anomaly Detection 3}

    \begin{figure}
        \centering
        \includegraphics[width=.99\textwidth]{img/ts/Anom1.png}
      
        \label{ANOM1}
    \end{figure}
    
    \begin{figure}
        \centering
        \includegraphics[width=.99\textwidth]{img/ts/Anom2.png}
      
        \label{ANOM2}
    \end{figure}
    
    
        
\end{frame}


%%%%%%%%%%%%%%%%%%%%%%%%%%%%%%%%%%%%%%%%%%%%%%%%%%%%%%%%%%%%%%%%%%%%%%%%%%


\begin{frame}{New Score}


    \begin{exampleblock}{\textsc{New\_Score = }}
        \begin{flalign*}
            n\_incidents\_week + n\_participants\_week
        \end{flalign*}
        \begin{flalign*}
            &n\_incidents\_week = \text{number of incidents in a specific week} \\
            &n\_participants\_week = \text{number of participants of all the incidents }\\
            &\text{in a specific week }
        \end{flalign*}
    \end{exampleblock}
    
\end{frame}
    

%%%%%%%%%%%%%%%%%%%%%%%%%%%%%%%%%%%%%%%%%%%%%%%%%%%%%%%%%%%%%%%%%%%%%%%%%%

\begin{frame}{ROC and Precision curve}
    \begin{figure}
        \centering
        \includegraphics[width=.99\textwidth]{img/ts/roc_pr.png}
    
        \caption{ROC curve and Precision curve of the classifier}
        \label{ROC}
    \end{figure}
\end{frame}


%%%%%%%%%%%%%%%%%%%%%%%%%%%%%%%%%%%%%%%%%%%%%%%%%%%%%%%%%%%%%%%%%%%%%%%%%%

\begin{frame}{Confusion Matrix}
    \begin{figure}
        \centering
        \includegraphics[width=.99\textwidth]{img/ts/confusion_matrix.png}
    
        \caption{Train and Test set Confusion Matrices}
        \label{MTX}
    \end{figure}
\end{frame}
    
    


\end{document}
